% =====================================================================
\section{The rate-independent framework}\label{sec:framework}
% =====================================================================

\subsection{The defining triple}\label{sec:ingredients}

Following~\cite{MielkeRoubicek}, a rate-independent system is specified
by a triple $(\Q, E, \Psi)$.

\begin{enumerate}[label=(\roman*),leftmargin=2.2em,itemsep=3pt]
  \item \textbf{State space.} The set of \emph{admissible states} is a
  closed and convex subset of a Banach space; throughout we take
  \begin{equation}\label{eq:state_space}
    \Q \subseteq \R^n .
  \end{equation}

  \item \textbf{Loading.} The interaction of the system with its
  environment is encoded by a prescribed, time-dependent
  \emph{loading path}
  \begin{equation}
    S \in W^{1,1}(0,T;\R^{+}) .
  \end{equation}

  \item \textbf{Stored energy.} The \emph{stored-energy functional}
  $E\colon \Q\times\R^{+}\to\R$ is assumed to admit the additive
  decomposition
  \begin{equation}\label{eq:energy}
    E(\bs q, S) \;=\; W(\bs q)\;-\;\bs q\cdot\bs\ell(S),
  \end{equation}
  where $W\colon\Q\to\R$ is the \emph{configurational free energy} and
  $\bs\ell\colon\R^{+}\to\R^n$ the \emph{interaction (loading)
  potential}.

  \item \textbf{Dissipation potential.} The \emph{dissipation
  potential} is a convex, lower semicontinuous map
  $\Psi\colon\R^n\to\R$ satisfying
  \begin{itemize}[leftmargin=1.4em,itemsep=1pt]
    \item \emph{normalisation}: $\Psi(\bs 0)=0$;
    \item \emph{positiveness}: $\Psi(\bs v)\ge 0$ for all $\bs v\in\R^n$;
    \item \emph{$1$-homogeneity}:
    $\Psi(\lambda\bs v)=\lambda\,\Psi(\bs v)$ for all $\lambda>0$.
  \end{itemize}
  These are the structural assumptions required
  in~\cite{Mielke2005,MielkeRoubicek} for the dissipation potential of
  a rate-independent system; the $1$-homogeneity is the property that
  distinguishes rate-independent from viscous (rate-dependent)
  evolution. A map satisfying only the first two conditions is called
  a dissipation \emph{pseudo-potential}.
\end{enumerate}

The forthcoming well-posedness results require, in addition, the
quantitative hypotheses of \cref{ass:standing}, which are collected in
\cref{sec:math}. The benchmark of \cref{sec:experiments} satisfies all
of them.

\subsection{Derived canonical quantities}\label{sec:derived}

From the triple $(\Q,E,\Psi)$ the following objects are canonically
constructed.

\paragraph{(a) Thermodynamic driving force.}
When $E$ is differentiable in $\bs q$, the negative energy gradient
defines the \emph{generalised thermodynamic force}
\begin{equation}\label{eq:force}
  \bs f(\bs q, S) \;:=\; -\frac{\partial E}{\partial \bs q}(\bs q, S)
  \;=\; \bs\ell(S) - \bs g(\bs q),
  \qquad
  \bs g(\bs q) := \frac{\partial W}{\partial\bs q}(\bs q).
\end{equation}
The force $\bs f$ is the variable conjugate to $\bs q$: it quantifies
the thermodynamic drive toward evolution of the state.

\paragraph{(b) Elastic domain.}
The subdifferential of $\Psi$ at the origin,
\begin{equation}\label{eq:elastic_domain}
  \E := \partial\Psi(\bs 0),
\end{equation}
is the closed convex set of forces the system can sustain without
evolving. It is the \emph{elastic domain} in force space; by convex
duality~\cite{rockafellar1970,ekeland1999convex} $\Psi$ is its support
function, $\Psi(\bs v)=\sup_{\bs f\in\E}\bs f\cdot\bs v$.

\paragraph{(c) Dissipation distance.}
The \emph{dissipation distance} is the minimal cost of a transition
between two states,
\begin{equation}\label{eq:diss_dist}
  \Dr(\bs q_1,\bs q_2)
  := \inf\Bigl\{\textstyle\int_0^1 \Psi(\dot\gamma(s))\dd s :
  \gamma\in\mathrm{AC}([0,1];\Q),\ \gamma(0)=\bs q_1,\ \gamma(1)=\bs q_2\Bigr\}.
\end{equation}

\paragraph{(d) Total dissipation.}
For a trajectory $\bs q\colon[0,T]\to\Q$ of bounded variation, the
\emph{total dissipation} is
\begin{equation}\label{eq:total_diss}
  \Diss_\Psi(\bs q;[0,t])
  := \sup\Bigl\{\textstyle\sum_{j=1}^{N}\Dr\bigl(\bs q(t_{j-1}),\bs q(t_j)\bigr):
  0=t_0<\dots<t_N=t,\ N\in\N\Bigr\}.
\end{equation}

\paragraph{(e) Power and work of the loading.}
The rate of energy change due to the loading at frozen state is
\begin{equation}\label{eq:power}
  \mathcal P_S(t) := \partial_t E(\bs q, S(t))
  = \partial_S E(\bs q, S)\,\dot S(t)
  = -\bs q\cdot\bs\ell'(S(t))\,\dot S(t),
\end{equation}
and the associated \emph{work} performed on the system over $[0,t]$ is
$\mathcal W_S(t):=\int_0^t\partial_S E(\bs q(s),S(s))\,\dot S(s)\dd s$.

\subsection{Governing equations: energetic formulation}
\label{sec:energetic}

We postulate the evolution through the following variational notion,
due to Mielke and Theil~\cite{MielkeTheil}.

\begin{definition}[Energetic solution]\label{def:energetic}
A function $\bs q\colon[0,T]\to\Q$ is an \emph{energetic solution} of
$(\Q,E,\Psi)$ for the loading $S$ if it satisfies:

\smallskip
\noindent\textbf{(S) Global stability.} For every $t\in[0,T]$ and every
$\tilde{\bs q}\in\Q$,
\begin{equation}\label{eq:stability}
  E(\bs q(t),S(t)) \;\le\; E(\tilde{\bs q},S(t)) + \Dr(\bs q(t),\tilde{\bs q}).
\end{equation}

\noindent\textbf{(E) Energy balance.} For every $t\in[0,T]$,
\begin{equation}\label{eq:energy_balance}
  E(\bs q(t),S(t)) + \Diss_\Psi(\bs q;[0,t])
  = E(\bs q(0),S(0)) + \int_0^t \partial_S E(\bs q(s),S(s))\,\dot S(s)\dd s.
\end{equation}
\end{definition}

Condition~(S) admits a variational reading: $\bs q(t)$ is a global
minimiser of the perturbed energy $\bs r\mapsto E(\bs r,S(t))+\Dr(\bs
q(t),\bs r)$. Condition~(E) states that, along the trajectory, the sum
of stored energy and cumulative dissipation equals the initial energy
plus the work supplied by the loading; it is an exact energy
conservation statement, examined thermodynamically in
\cref{sec:thermo}.

\subsection{Governing equations: subdifferential formulation}
\label{sec:subdiff}

Since $\Psi$ is convex but generally non-smooth at $\dot{\bs q}=\bs 0$,
the pointwise force balance takes the form of a differential inclusion.

\begin{proposition}[Biot inclusion]\label{prop:biot}
Let $\bs q\in\mathrm{AC}([0,T];\Q)$ satisfy \textnormal{(S)}
and~\textnormal{(E)}, with $E(\cdot,S)$ differentiable. Then, for a.e.\
$t\in[0,T]$,
\begin{equation}\label{eq:biot}
  \bs 0 \;\in\; \partial\Psi(\dot{\bs q}(t)) + \frac{\partial E}{\partial\bs q}(\bs q(t),S(t)),
\end{equation}
equivalently $\bs f(\bs q,S)\in\partial\Psi(\dot{\bs q})$. Conversely,
if $E(\cdot,S)$ is convex, then~\eqref{eq:biot} together with the
local stability $\bs f(\bs q(t),S(t))\in\E$ implies
\textnormal{(S)}--\textnormal{(E)}.
\end{proposition}

Equation~\eqref{eq:biot} is the (generalised) Biot
equation~\cite{Biot1956}; for $\Psi=\rho\|\cdot\|$ it coincides with the
sweeping process of Moreau~\cite{moreau1977} and, in one dimension, with
the play operator of hysteresis~\cite{Krasnoselskii,visintin1994}. Its
analysis rests on the \emph{saturation identity}, the algebraic
signature of $1$-homogeneity: if $\bs f\in\partial\Psi(\bs v)$ then
$\bs f\cdot\bs v=\Psi(\bs v)$ (\cref{lem:saturation} in
\cref{app:framework}). Combined with \cref{prop:biot}, it yields, along
any solution,
\begin{equation}\label{eq:saturation_solution}
  \bs f(\bs q,S)\cdot\dot{\bs q}=\Psi(\dot{\bs q}),
\end{equation}
i.e.\ the dissipation is exactly balanced by the energy release at
frozen loading. This is the cornerstone of the thermodynamic
interpretation of \cref{sec:thermo}. The proof of \cref{prop:biot} is
given in \cref{app:framework}.
